% Part: first-order-logic
% Chapter: proof-systems
% Section: tableaux

\documentclass[../../../include/open-logic-section]{subfiles}

\begin{document}

\iftag{FOL}
      {\olfileid{fol}{prf}{tab}}
      {\olfileid{pl}{prf}{tab}}

\olsection{\usetoken{P}{tableau}}

Nalika akeh sistem !!{derivation} nggarap susunan !!{sentence}s,
!!{tableau}s nggarap !!{signed formula}s. !!^a{signed formula} yaiku
pasangan kang dumadi saka tandha nilai kabeneran ($\True$ utawa
$\False$) lan !!a{sentence}
\[
\sFmla{\True}{!A} \text{ utawa } \sFmla{\False}{!A}.
\]
!!^a{tableau} dumadi saka !!{signed formula}s kang disusun dadi wit
kang nyabang mudhun. Tableau diwiwiti nganggo sawetara
\emph{asumsi}, banjur diterusake dening !!{signed formula}s kang
diasilake saka salah siji !!{signed formula}s ing ndhuwure kanthi
ngetrapake salah siji aturan inferensi. Saben aturan ngidini kita
nambahake siji utawa luwih !!{signed formula}s ing pucuk sawijining
cabang, utawa rong !!{signed formula}s jejer---ing kasus iki, cabang
pecah dadi loro, lan rong !!{signed
  formula}s kang ditambahake dadi pucuk rong cabang kasebut.

Aturan kang ditrapake marang !!{signed formula} kompleks ngasilake
tambahan !!{signed formula}s kang dadi sub-!!{formula}s langsung.
Aturan kasebut teka minangka pasangan, siji kanggo saben tandha.
Tuladhane, aturan $\TRule{\True}{\land}$ ditrapake marang
$\sFmla{\True}{!A \land !B}$ lan ngidini loro !!{signed formula}s,
$\sFmla{\True}{!A}$ lan~$\sFmla{\True}{!B}$, ditambahake ing pucuk
saben cabang kang ngemot $\sFmla{\True}{!A \land !B}$, dene aturan
$\TRule{\False}{\land}$ ngidini sawijining cabang dipecah kanthi
nambahake $\sFmla{\False}{!A}$ lan $\sFmla{\False}{!B}$ kanthi jejer.
Cathetan koreksi sumber OLPL-002: jeneng aturan ing sumber Inggris
ngemot formula wutuh; jeneng iki dibalekake dadi aturan konjungsi
mawa tandha salah.
!!^a{tableau} katutup yen saben cabange ngemot pasangan
!!{signed formula}s kang cocog, yaiku $\sFmla{\True}{!A}$ lan
$\sFmla{\False}{!A}$.

Relasi $\Proves$ adhedhasar !!{tableau}s ditemtokake kaya mangkene:
$\Gamma \Proves !A$ yen lan mung yen ana himpunan winates
$\Gamma_0 = \{!B_1, \dots, !B_n\} \subseteq \Gamma$ saengga ana
!!{tableau} katutup kanggo asumsi
\[
\{\sFmla{\False}{!A}, \sFmla{\True}{!B_1}, \dots, \sFmla{\True}{!B_n}\} 
\]
Tuladhane, iki !!{tableau} katutup kang nuduhake yen $\Proves
(!A \land !B) \lif !A$:
\begin{oltableau}
  [\sFmla{\False}{(\formula{A} \land \formula{B}) \lif \formula{A}}, just = \TAss
    [\sFmla{\True}{\formula{A} \land \formula{B}}, just = {\TRule{\False}{\lif}[1]}
      [\sFmla{\False}{\formula{A}}, just = {\TRule{\False}{\lif}[1]}
        [\sFmla{\True}{\formula{A}}, just = {\TRule{\True}{\land}[2]}
          [\sFmla{\True}{\formula{B}}, just = {\TRule{\True}{\land}[2]}, close]
        ]
      ]
    ]
  ]
\end{oltableau}
Cathetan koreksi sumber OLPL-003: rong label ekspansi konjungsi ing
sumber Inggris nyebut implikasi mawa tandha bener; label kasebut
dibalekake dadi konjungsi mawa tandha bener, cocog karo formula ing
larik induke.

Himpunan $\Gamma$ ora konsisten ing kalkulus !!{tableau} yen ana
!!{tableau} katutup kanggo asumsi
\[
\{\sFmla{\True}{!B_1}, \dots, \sFmla{\True}{!B_n}\} 
\]
kanggo sawijining $!B_i \in \Gamma$.

!!^{tableau}s ditemokake ing taun 1950-an kanthi mandhiri dening Evert
Beth lan Jaakko Hintikka, banjur disederhanakake lan dipopulerake dening
Raymond Smullyan. Tableau gampang banget digunakake, amarga mbangun
!!a{tableau} iku prosedur kang sistematis banget. Amarga sipate kang
sistematis, !!{tableau}s uga gampang ditindakake nganggo komputer.
Nanging, !!a{tableau} kerep angel diwaca lan gegayutane karo bukti
kadhangkala ora gampang katon. Pendekatan iki uga cukup umum, lan akeh
logika duwe sistem !!{tableau}. !!^{tableau}s uga mbiyantu kita nemokake
!!{structure}s kang nyukupi (himpunan) !!{sentence}s kang diwenehake:
yen himpunan iku bisa dicukupi, himpunan iku ora bakal duwe
!!{tableau} katutup, yaiku saben !!{tableau} bakal duwe cabang kang
bukak. !!{structure} kang nyukupi bisa ``diwaca'' saka cabang bukak,
angger saben aturan kang bisa ditrapake wis ditrapake ing cabang kasebut.
Uga ana gegayutan kang raket banget karo kalkulus sekuen: pokoke,
!!{tableau} katutup iku !!{derivation} kalkulus sekuen kang dipadhetake
lan ditulis kuwalik saka ndhuwur menyang ngisor.

\end{document}
